\documentclass[10pt]{article}
\usepackage[T1]{fontenc}
\usepackage[utf8]{inputenc}
\usepackage[french]{babel}
\usepackage{amsmath,amssymb,mathrsfs,mathpazo}
\usepackage{geometry}
\geometry{a4paper,margin=2cm}
\usepackage{graphicx,array,tabularx,tabularray,colortbl,multirow,multicol,booktabs,enumitem,pifont,ulem,fancybox,diagbox,cancel}
\usepackage{pstricks,pstricks-add,pst-plot,pst-tree,pst-node,pst-text,pst-func,pst-math,pst-bspline,pst-3dplot,pst-eucl,pst-solides3d}
\let\pscalculate\undefined
\usepackage{pst-bezier}
\usepackage{tikz,pgfplots,tkz-tab}
\usetikzlibrary{arrows,arrows.meta,patterns,calc,positioning,trees,decorations.markings,shapes.geometric,shapes.misc,shapes.arrows,intersections,decorations.pathmorphing,decorations.pathreplacing,backgrounds,shadows,babel}
\pgfplotsset{compat=1.18}
\usepgfplotslibrary{fillbetween}
\usepackage[np]{numprint}
\usepackage{listings,makecell,pas-tableur,twoopt}
\let\barre\undefined
\usepackage[pstricks]{tablvar}
\definecolor{darkgreen}{RGB}{0,100,0}
\providecommand{\red}{\color{red}}
\providecommand{\blue}{\color{blue}}
\providecommand{\green}{\color{green}}
\providecommand{\black}{\color{black}}
\providecommand{\white}{\color{white}}
\providecommand{\gray}{\color{gray}}
\providecommand{\cyan}{\color{cyan}}
\providecommand{\magenta}{\color{magenta}}
\providecommand{\yellow}{\color{yellow}}
\providecommand{\violet}{\color{violet}}
\providecommand{\orange}{\color{orange}}
\providecommand{\euro}{\texteuro}
\providecommand{\EUR}{\texteuro}
\providecommand{\ieme}{\textsuperscript{e}}
\providecommand{\degres}{\ensuremath{^\circ}}
\providecommand{\degree}{\ensuremath{^\circ}}
\providecommand{\vv}[1]{\overrightarrow{#1}}
\providecommand{\Coord}[1]{#1}
\providecommand{\decosix}{}
\providecommand{\starredbullet}{\ensuremath{\bullet}}
\pagestyle{empty}
\setlength{\parindent}{0pt}
\usepackage{adjustbox}\DeclareUnicodeCharacter{2236}{:}\DeclareUnicodeCharacter{2190}{\ensuremath{\leftarrow}}
\providecommand{\R}{\mathbb{R}}
\providecommand{\N}{\mathbb{N}}
\providecommand{\Z}{\mathbb{Z}}
\providecommand{\Q}{\mathbb{Q}}
\providecommand{\D}{\mathbb{D}}
\providecommand{\C}{\mathbb{C}}
\providecommand{\e}{{\mathrm{e}}}
\providecommand{\vect}[1]{\overrightarrow{#1}}
\providecommand{\vectt}[1]{\overrightarrow{\mathrm{#1}}}
\providecommand{\vvt}[1]{\overrightarrow{\mathrm{#1}}}
\providecommand{\np}[2][]{{#2~#1}}
\providecommand{\vv}[1]{{\overrightarrow{#1}}}
\renewcommand{\d}{{\mathrm{d}}}
\providecommand{\degres}{{^\circ}}
\providecommand{\degre}{{^\circ}}
\providecommand{\blue}{\color{blue}}
\providecommand{\red}{\color{red}}
\providecommand{\enskip}{\quad}
\providecommand{\llbracket}{\lbrack\!\lbrack}
\providecommand{\rrbracket}{\rbrack\!\rbrack}
\providecommand{\euro}{{\text{€}}}
\providecommand{\no}{{\text{n°}}}
\providecommand{\interval}[2]{{[#1;#2]}}
\providecommand{\intervFO}[2]{{[#1;#2[}}
\providecommand{\dotfill}{\ldots}

\providecommand{\ds}{\displaystyle}
\providecommand{\textding}[1]{#1}
\providecommand{\Oij}{\(\mathrm{O};\vec{\imath},\vec{\jmath}\)}
\providecommand{\Oijk}{\(\mathrm{O};\vec{\imath},\vec{\jmath},\vec{k}\)}
\providecommand{\Ouv}{\(\mathrm{O};\vec{u},\vec{v}\)}
\providecommand{\intervFF}[2]{\ensuremath{[#1;#2]}}
\renewcommand{\euro}{\ifmmode\text{\texteuro}\else\texteuro\fi}
\DeclareUnicodeCharacter{00B0}{\ifmmode^{\circ}\else\textdegree\fi}
\DeclareUnicodeCharacter{20AC}{\ifmmode\text{\texteuro}\else\texteuro\fi}
% Énoncés : Éducation nationale. Transcription et corrigés : APMEP, Denis Vergès et contributeurs. Découpage et classement : CoopMaths, Sébastien Lozano, Éric Elter et contributeurs. Adaptation de la présentation : Maths Explorer.
% CC BY-SA : https://coopmaths.fr/www/about/
% https://coopmaths.fr/alea/static/bac/2026/tex/bac_2026_05_sujet2_ameriquenord_1_cor.tex
\begin{document}
\noindent\textit{APMEP / CoopMaths. CC BY-SA. Présentation adaptée par Maths Explorer.}\par\medskip
Un supermarché dispose d'un stock de tomates provenant de deux fournisseurs A et B.

II a été constaté que :

\begin{itemize}[label=$\bullet~$]
\item 91\,\% du stock de tomates est commercialisable ;
\item 60\,\% du stock de tomates provient du fournisseur A ;
\item parmi les tomates provenant du fournisseur A, la proportion de tomates commercialisables
est de 95\,\%.
\end{itemize}

On choisit au hasard une tomate dans le stock.

On désigne par :
\begin{itemize}[label=$\bullet~$]
\item $A$ l'évènement \og La tomate provient du fournisseur A \fg ;
\item $B$ l'évènement \og La tomate provient du fournisseur B \fg ;
\item $C$ l'évènement \og La tomate est commercialisable \fg.
\end{itemize}

Pour un évènement quelconque $E$, on note $P(E)$ la probabilité de $E$.

\medskip

\textbf{Partie A}

\medskip

\begin{enumerate}
\item On complète l'arbre ci-dessous.


\begin{center}
\pstree[treemode=R,levelsep=3.5cm, treesep=0.7cm,nodesepA=0pt,nodesepB=4pt,nrot=:U]{\TR{}}
{\pstree[nodesepA=4pt]{\TR{$A$}\naput{\blue$0,60$}}
	{\TR{$C$} \naput{\blue$0,95$}
	\TR{$\overline{C}$} \nbput{\blue$1-0,95=0,05$}
	}
\pstree[nodesepA=4pt]{\TR{$B$}\nbput{\blue$1-0,60=0,40$}}
	{\TR{$C$} 
	\TR{$\overline{C}$} 
	}
}
\end{center}























\item
	\begin{enumerate}
		\item La probabilité que la tomate choisie soit commercialisable et provienne du fournisseur A est
$P\left (A\cap C\right ) = P\left (A\right ) \times P_{A}\left (C\right ) = 0,60\times 0,95=0,57$.
		
		\item On cherche  $P_B\left (C\right )$.
D'après le texte: $P\left (C\right ) = 0,91$.

D'après la formule des probabilités totales, on a:

$P\left (C\right )=P\left (A\cap C\right )+P\left (B\cap C\right )$ donc $P\left (C\right )= P\left (A\cap C\right )+P\left (B\right )\times P_{B}\left (C\right )$.

Donc: $0,91=0,57+0,40\times P_{B}\left (C\right )$, 
soit $0,34=0,40\times P_{B}\left (C\right )$,
et donc: $P_{B}\left (C\right ) = \dfrac{0,34}{0,40}=0,85$.
		
		\item La tomate choisie est non commercialisable. Le responsable des achats estime qu'il y a deux fois moins de chance qu'elle provienne du fournisseur A que du fournisseur B.
		
Il faut donc comparer $P_{\overline{C}}\left (A\right )$ et  $P_{\overline{C}}\left (B\right )$.

$P_{\overline{C}}\left (A\right ) = \dfrac{P\left ( A \cap \overline{C}\right )}{P\left (\overline{C}\right )}
= \dfrac{0,60\times 0,05}{1-0,91} = \dfrac{0,03}{0,09}= \dfrac{1}{3}$ 

$P_{\overline{C}}\left (B\right ) = \dfrac{P\left ( B \cap \overline{C}\right )}{P\left (\overline{C}\right )}
= \dfrac{0,40\times \left (1-0,85\right )}{0,09} = \dfrac{0,06}{0,09}= \dfrac{2}{3}$ 

$P_{\overline{C}}\left (A\right ) = \dfrac{1}{2}\times P_{\overline{C}}\left (B\right )$ donc le responsable commercial a raison.

	\end{enumerate}
\end{enumerate}

\bigskip

\textbf{Partie B}

\medskip

On rappelle que 9\,\% des tomates du stock ne sont pas commercialisables.

\medskip

\begin{enumerate}
\item On prend $15$ tomates dans le stock au hasard et de manière indépendante. On considère que le stock est suffisamment important pour qu'on puisse assimiler ce prélèvement à un tirage aléatoire avec remise.
On note $X$ la variable aléatoire égale au nombre de tomates non commercialisables dans cet échantillon de $15$~tomates.

	\begin{enumerate}
		\item On admet que la variable aléatoire $X$ suit une loi binomiale; ses paramètres sont $n=15$ et $p=0,09$.
		
		\item La probabilité qu'exactement deux tomates soient non commercialisables est:
		
		$P\left (X=2\right ) = \ds\binom{15}{2}\times 0,09^{2}\times \left (1-0,09\right )^{15-2}$ qui a pour arrondi au millième $0,250$.



		\item La probabilité qu'au plus deux tomates soient non commercialisables est:
		
		$P\left (X \leqslant 2\right ) \approx 0,853$ (à la calculatrice).


	\end{enumerate}
\item On constitue désormais un échantillon de $n$ tomates, toujours dans les mêmes conditions, où $n$ désigne un entier naturel non nul.

On note $X_n$, la variable aléatoire égale au nombre de tomates non commercialisables et $F_n$ la variable aléatoire égale à la fréquence de tomates non commercialisables dans cet échantillon de $n$ tomates.

On a donc $F_n = \dfrac{X_n}{n}$.

On admet que la variable aléatoire $X_n$, suit la loi binomiale de paramètres $n$ et $0,09$.

	\begin{enumerate}
		\item 
		\begin{list}{\textbullet}{}
		\item Puisque $X_n$ suit la loi binomiale de paramètres $n$ et $0,09$, alors l'espérance de $X_n$ est, par propriété : $E(X_n) = np = 0,09n$.

On a ensuite $F_n = \dfrac{1}{n} \times X_n$,

donc, par linéarité de l'espérance : $E(F_n) = \dfrac{1}{n} E(X_n) = \dfrac{1}{n}\times 0,09n = 0,09$.

\item Puisque $X_n$ suit la loi binomiale de paramètres $n$ et $0,09$, alors la variance de $X_n$ est, par propriété : $V(X_n) = np(1-p) = 0,09\times (1-0,09)\times n = 0,0819n$.

D'après la formule: $V(aX) = a^2V(X)$, on a:
$V\left (F_n\right ) = V\left ( \dfrac{X_n}{n} \right ) = V\left (\dfrac{1}{n} X_n\right ) = \dfrac{V\left (X_n\right )}{n^2}$.

Donc $V\left (F_n\right )= \dfrac{V\left (X_n\right )}{n^2}	= \dfrac{0,0819n}{n^2}=\dfrac{0,0819}{n}$.
		\end{list}

		\item On veut démontrer que $P\left (0,04 < F_n < 0,14\right ) \geqslant 1 - \dfrac{32,76}{n}$.

$\begin{aligned}
0,04 < F_n < 0,14
& \iff 0,09-0,05 < F_n < 0,09+0,05
\iff -0,05 < F_n-0,09 < 0,05\\
& \iff \left | F_n-0,09 \right | < 0,05
\end{aligned}$

donc 
$P\left (0,04 < F_n < 0,14\right ) = P\left (\left | F_n-0,09 \right | < 0,05\right ) = 1-P \left (\left | F_n-0,09 \right |\geqslant 0,05\right ) $

La variable aléatoire $F_n$ a pour espérance $E\left (F_n\right )$ et pour variance $V\left (F_n\right )$ donc, d'après l'inégalité de Bienaymé-Tchebychev, on a:

pour tout $\delta\in\, \left ]0\;;\;+\infty\strut \right [$, 
$P\left (\left |F_n-E\left (F_n\right)\strut \right | \geqslant \delta \right ) \leqslant \dfrac{V\left (F_n\right  )}{\delta^2}$.

$E\left (F_n\right )=0,09$ et $V\left (F_n\right )= \dfrac{0,0819}{n}$ donc  l'inégalité de Bienaymé-Tchebychev devient:

pour tout $\delta\in\, \left ]0\;;\;+\infty\strut \right [$, 
$P\left ( \left |F_n-0,09\strut \right | \geqslant \delta \right ) \leqslant \dfrac{\frac{0,0819}{n}}{\delta^2}$.

Pour $\delta=0,05$, l'inégalité:
$P\left ( \left |F_n-0,09\strut \right | \geqslant 0,05 \right ) \leqslant \dfrac{\frac{0,0819}{n}}{0,05^2}$,
soit
$P\left ( \left |F_n-0,09\strut \right | \geqslant 0,05 \right ) \leqslant \dfrac{32,76}{n}$		
		
$P\left (0,04 < F_n < 0,14\right ) =  1-P \left (\left | F_n-0,09 \right | \geqslant 0,05\right )$	

$P\left ( \left |F_n-0,09\strut \right | \geqslant 0,05 \right ) \leqslant \dfrac{32,76}{n}$
donc
$-P\left ( \left |F_n-0,09\strut \right | \geqslant 0,05 \right ) \geqslant -\dfrac{32,76}{n}$

donc
$1-P\left ( \left |F_n-0,09\strut \right | \geqslant 0,05 \right ) \geqslant 1 -\dfrac{32,76}{n}$
et donc
$P\left (0,04 < F_n < 0,14\right ) \geqslant 1 -\dfrac{32,76}{n}$

		\item Le responsable des achats prélève dans le stock un échantillon de $500$ tomates. Il s'aperçoit que $55$ tomates ne sont pas commercialisables.



\begin{list}{\textbullet}{Pour $n=500$:}
\item La fréquence de tomates non commercialisables est $F=\dfrac{55}{500}=0,11$; on remarque que $0,04 \leqslant F \leqslant 0,14$
\item $1-\dfrac{32,76}{n}= 1- \dfrac{32,76}{500} = \np{0,93448}$
\end{list}

On a donc $P\left (0,04 \leqslant F_{500} \leqslant 0,14\right ) \geqslant \np{0,93448}$.

Il y a donc une probabilité supérieure à $93,448$\,\% que la fréquence de tomates non commercialisables soit comprise entre $0,04$ et $0,14$.

Comme $0,11$ est compris  entre $0,04$ et $0,14$, le résultat est conforme à ce qu'attendait le responsable.

	\end{enumerate}
\end{enumerate}

\medskip

\end{document}
